\section{Step 26 Protect inputs and inspect actual stores}\label{maze-s26}

Final byte equality cannot reveal a forbidden write that is later restored. The checks also observe execution stores to maze inputs and guard words.

\subsection{Write or inspect}

In a temporary copy of your program, write a wall word and restore its old value before continuing the search. Run the public checks and compare the failure with a plain final-memory assertion. Remove the extra stores afterward.

\subsection{Expected observation}

The input-store trace rejects the mutation even when the final maze is unchanged. Register scratch choices and helper tables remain implementation details.

\subsection{Check your understanding}

Why observe stores as well as final bytes?

\begin{itemize}
  \item \textbf{A} To force one register allocation.
  \item \textbf{B} To detect a forbidden write followed by restoration.
  \item \textbf{C} To choose a unique shortest path.
\end{itemize}
